\subsection{Compact Operators}\label{sec:compact-operators}\label{sec:integral-operators}

\begin{Definition}
\label{def:compact-operator}
Let $X$ and $Y$ be Banach spaces, and let $(\mathcal{D}(A),A)$ be a linear operator from $X$ to $Y$. Then we say the operator $A$ is \textbf{compact} if it maps bounded sets into precompact sets; i.e., if for every bounded set $\Omega\subset\mathcal{D}(A)$, we have $\overline{A(\Omega)}\subset Y$ compact.
\end{Definition}

It is often convenient to characterize compact operators in terms of sequences rather than in terms of sets.

\begin{Theorem}
\label{thm:compact-operator-sequential-characterization}
An operator $(\mathcal{D}(A),A)$ from $X$ to $Y$ is compact if and only if it is sequentially compact; i.e., if and only if given any bounded sequence $x_n$ in $\mathcal{D}(A)$, it follows that $A(x_n)$ has a convergent subsequence.
\end{Theorem}

\begingroup
\postdisplaypenalty=10000
\renewcommand{\endProof}{\par\nobreak\hfill$\blacksquare$\par\vspace{0.1cm}}
\begin{Proof}
The proof of this theorem follows directly from the topological result that a precompact set can be characterized by sequences; i.e., a set $S$ in a normed linear space is precompact if and only if every sequence contained in $S$ has a convergent subsequence.
\end{Proof}
\endgroup

As we shall see below, the most fundamental examples of compact operators are integral operators. However, we shall need to develop a bit of machinery in order to study them more fully. In the meantime, we have been provided with some very important examples of compact operators by our study of compact imbeddings in Section~\ref{subsec:sobolev-compactness-properties}. In order to interpret them we need the following lemma.

\begin{Lemma}
\label{lem:compact-embedding-identity-operator}
Let $X$ and $Y$ be Banach spaces. Then $X$ is compactly imbedded in $Y$ if and only if the identity mapping from $X$ to $Y$ is well defined and compact.
\end{Lemma}

The proof follows immediately from Definition~\ref{def:compact-imbedding} and the definition of the identity mapping in Example~\ref{ex:identity-operator-and-inclusion}.

\begin{Example}
\label{ex:sobolev-continuous-target-compact-operator}
It follows from Theorem~\ref{thm:sobolev-continuous-target-compact-embedding} that if $k>m/2$ and $\Omega\subset\mathbb{R}^m$ is bounded and has the $k$-extension property, then the identity mapping from $H^k(\Omega)$ to $C_b(\overline\Omega)$ is compact. Thus by Theorem~\ref{thm:compact-operator-sequential-characterization}, every sequence of functions $u_n$ that is bounded in the $H^k(\Omega)$ norm has a uniformly convergent subsequence.
\end{Example}

\begin{Example}
\label{ex:successive-sobolev-orders-compact-operator}
It follows from Theorem~\ref{thm:compact-successive-sobolev-orders} that if $k$ is a non-negative integer and $\Omega\subset\mathbb{R}^m$ is bounded and has the $k+1$-extension property, then the identity mapping from $H^{k+1}(\Omega)$ to $H^k(\Omega)$ is compact. Using Theorem~\ref{thm:compact-operator-sequential-characterization} again, we see that every sequence of functions $u_n$ that is bounded in the $H^{k+1}(\Omega)$ norm has a subsequence that converges strongly in the $H^k(\Omega)$ norm.
\end{Example}

We now obtain the following elementary result.

\begin{Lemma}
\label{lem:compact-operator-boundedness}
Every compact operator is bounded.
\end{Lemma}

\begingroup
\postdisplaypenalty=10000
\renewcommand{\endProof}{\par\nobreak\hfill$\blacksquare$\par\vspace{0.1cm}}
\begin{Proof}
Suppose not, then there is a sequence $x_n\in\mathcal{D}(A)$ such that $\|x_n\|=1$ and $\|A(x_n)\|\to\infty$. In fact, by eliminating superfluous elements of the sequence and relabeling, we can ensure that $\|A(x_{n+1})\|>\|A(x_n)\|+1$. Thus, no subsequence of $A(x_n)$ could converge since no subsequence could be Cauchy.
\end{Proof}
\endgroup

Recall that by Theorem~\ref{thm:bounded-operator-norm-preserving-extension}, every bounded operator can be extended to all of $X$ without changing its norm. We leave it to the reader to show that when a compact operator is extended using the methods described in the proof of Theorem~\ref{thm:bounded-operator-norm-preserving-extension}, the extended operator is also compact (Problem~\ref{ex:compact-operator-norm-preserving-extension}). Thus, we will usually assume that a compact operator is in $\mathcal{L}(X,Y)$.

Note that Lemma~\ref{lem:compact-operator-boundedness} and Lemma~\ref{lem:linear-mapping-bounded-continuity-equivalence} tell us that every compact operator is continuous. However, the converse of this result is false. In particular, we have the following.

\begin{Lemma}
\label{lem:infinite-dimensional-identity-noncompact}
If $X$ is any infinite-dimensional Banach space, then the identity operator is not compact.
\end{Lemma}

\begingroup
\postdisplaypenalty=10000
\renewcommand{\endProof}{\par\nobreak\hfill$\blacksquare$\par\vspace{0.1cm}}
\begin{Proof}
The proof follows immediately from the fact that in an infinite-dimensional space, the unit ball is not compact. We prove this only in the case of an infinite-dimensional Hilbert space and leave the general result to the reader (Problem~\ref{ex:infinite-dimensional-identity-noncompact}). Recall that, by Corollary~\ref{cor:hilbert-basis-existence}, in an infinite-dimensional Hilbert space there exists an infinite orthonormal set $\{x_i\}_{i=1}^\infty$. This set is contained in the closed unit ball, and if $x_i$ and $x_j$ are two distinct elements of the basis, we have $\|x_i-x_j\|^2=2$. Thus, no subsequence of $x_i$ could converge since no subsequence could be Cauchy.
\end{Proof}
\endgroup
